\documentclass[notes,10pt, aspectratio=169]{beamer}

\usepackage{pgfpages}
% These slides also contain speaker notes. You can print just the slides,
% just the notes, or both, depending on the setting below. Comment out the want
% you want.
% Note visibility is selected after package loading at the end of this preamble.
%\setbeameroption{show only notes} % Only notes
%\setbeameroption{show notes on second screen=right} % Both

\usepackage{helvet}
\usepackage[default]{lato}
\usepackage{array}
\usepackage{eurosym}
\usepackage{setspace}
\usepackage{animate}
\usepackage{multicol}
\usepackage{breqn}
\usepackage{subcaption}
\usepackage{caption}
\captionsetup{font=scriptsize, justification=raggedright,singlelinecheck=false}
\usepackage[export]{adjustbox}
\usepackage[scale=2]{ccicons}
\usepackage{tikz}
\newenvironment{wideitemize}{\itemize\addtolength{\itemsep}{10pt}}{\enditemize}
\usepackage{natbib}

\usepackage{tikz}
\usepackage{verbatim}
\setbeamertemplate{note page}{\pagecolor{yellow!5}\insertnote}
\usepackage{amsmath}
%\usepackage{mathpazo}
\usepackage{hyperref}
\usepackage{lipsum}
\usepackage{multimedia}
\usepackage{multirow}
\usepackage{graphicx}
\usepackage{dcolumn}
\usepackage{bbm}
\newcolumntype{d}[0]{D{.}{.}{5}}

\usepackage{changepage}
\usepackage{appendixnumberbeamer}
% \newcommand{\beginbackup}{
%   \newcounter{framenumbervorappendix}
%   \setcounter{framenumbervorappendix}{\value{framenumber}}
%   \setbeamertemplate{footline}
%   {
%      \leavevmode%
%      \hline
%      box{%
%       \begin{beamercolorbox}[wd=\paperwidth,ht=2.25ex,dp=1ex,right]{footlinecolor}%
% %         \insertframenumber  \hspace*{2ex} 
%       \end{beamercolorbox}}%
%      \vskip0pt%
%   }
%  }

 \setbeamertemplate{footline}{
    \hbox{%
    \begin{beamercolorbox}[wd=\paperwidth,ht=3ex,dp=1.5ex,leftskip=2ex,rightskip=2ex]{page footer}%
        %\usebeamerfont{title in head/foot}%
        %\insertshorttitle \hfill
        %\insertsection \hfill
        %\tableofcontents[currentsection] \hfill
        \insertsectionnavigationhorizontal{0.7\paperwidth}{}{} \hfill 
        \insertframenumber{} / \inserttotalframenumber
    \end{beamercolorbox}}%
}
\newcommand{\backupend}{
   \addtocounter{framenumbervorappendix}{-\value{framenumber}}
   \addtocounter{framenumber}{\value{framenumbervorappendix}} 
}


\usepackage{graphicx}
\usepackage[space]{grffile}
\usepackage{booktabs}

% These are my colors -- there are many like them, but these ones are mine.
\definecolor{blue}{RGB}{0,114,178}
\definecolor{red}{RGB}{213,94,0}
\definecolor{yellow}{RGB}{240,228,66}
\definecolor{green}{RGB}{0,158,115}

\setbeamercolor{section in head/foot}{fg=blue}

\useinnertheme[shadow]{rounded}
\setbeamercolor{block title example}{fg=black,bg=orange!50!white}
\setbeamercolor{block body example}{fg=black, bg=orange!30!white}


\hypersetup{
  colorlinks=false,
  linkbordercolor = {white},
  linkcolor = {blue}
}


%% I use a beige off white for my background
\definecolor{MyBackground}{RGB}{255,253,218}

%% Uncomment this if you want to change the background color to something else
%\setbeamercolor{background canvas}{bg=MyBackground}

%% Change the bg color to adjust your transition slide background color!
\newenvironment{transitionframe}{
  \setbeamercolor{background canvas}{bg=yellow}
  \begin{frame}}{
    \end{frame}
}

\setbeamercolor{frametitle}{fg=blue}
\setbeamercolor{title}{fg=black}
%\setbeamertemplate{footline}[frame number]
\setbeamertemplate{navigation symbols}{} 
\setbeamertemplate{itemize items}{-}
\setbeamercolor{itemize item}{fg=blue}
\setbeamercolor{itemize subitem}{fg=blue}
\setbeamercolor{enumerate item}{fg=blue}
\setbeamercolor{enumerate subitem}{fg=blue}
\setbeamercolor{button}{bg=white,fg=blue,}



% If you like road maps, rather than having clutter at the top, have a roadmap show up at the end of each section 
% (and after your introduction)
% Uncomment this is if you want the roadmap!
% \AtBeginSection[]
% {
%    \begin{frame}
%        \frametitle{Roadmap of Talk}
%        \tableofcontents[currentsection]
%    \end{frame}
% }
\setbeamercolor{section in toc}{fg=blue}
\setbeamercolor{subsection in toc}{fg=red}
%\setbeamersize{text margin left=1em,text margin right=1em} 
\setbeamersize{text margin left=0.7cm,text margin right=0.7cm} 


\title[]{\textcolor{blue}{Seminar 2:}}
\subtitle{Classroom problems and worked solutions}
\author{25117 - Econometrics}
\institute[shortinst]{Universitat Pompeu Fabra}

%\date{September 25th, 2024}


% Each problem is followed by worked solution slides. Set \showsolutionsfalse
% below to compile a problem-only deck. Hidden \note entries retain instructor
% guidance. All data are hypothetical; no external assets are needed.
\newif\ifshowsolutions
\showsolutionstrue
\newcommand{\E}{\mathbb{E}}
\newcommand{\Var}{\operatorname{Var}}
\newcommand{\Cov}{\operatorname{Cov}}
\newcommand{\Corr}{\operatorname{Corr}}
\newcommand{\answerpart}[1]{\textcolor{blue}{\textbf{#1.}}\,}

\setbeameroption{hide notes} % Speaker notes hidden; worked solutions are visible.

\begin{document}

\begin{frame}[plain]
\maketitle
\end{frame}
\note{Instructor guide. This seminar follows Lecture2v2.tex: sample dispersion and standard errors, two-sided hypothesis tests, testing errors and power, confidence intervals, comparisons between independent groups, randomized experiments, and sample descriptive statistics. All numerical examples are invented for teaching. Each problem is followed by visible worked solutions and retains a hidden instructor note. The normal approximation is used for inference unless an exact normal model with known variance is explicitly supplied.}

\begin{frame}{The issue: evidence about student and school performance}
\begin{wideitemize}
\item We want to assess claims about average test scores and differences between groups.
\item Sampling uncertainty determines how precisely we can estimate a mean and how much evidence a difference provides.
\item For each problem, write the formula, show the calculation, and state what the result tells us about the population.
\end{wideitemize}
\begin{exampleblock}{Classroom problems}
Work through each problem before viewing its solution. Use the assumptions and normal probabilities supplied in the question.
\end{exampleblock}
\end{frame}
\note{Suggested modules: Problems 1--2 cover dispersion and precision; 3--5 cover tests and their errors; 6--7 cover confidence intervals; 8--9 compare groups; 10--12 apply sample summaries and inference. Every solution part is numbered to match its question.}

\section{Precision}

\begin{frame}{Problem 1: How dispersed are five test scores?}
Five students obtain the following scores, measured in points:
\[
60,\quad60,\quad70,\quad80,\quad80.
\]
Treat the scores as i.i.d. observations from a population with finite variance.
\begin{enumerate}
\setlength{\itemsep}{9pt}
\item Find the sample mean and each student's deviation from it. Check the sum of the deviations.
\item Calculate the sample variance using the divisor $n-1$, and then the sample standard deviation. State their units.
\item What variance would the divisor $n$ give? Explain why $n-1$ is used when estimating the population variance, and identify the remaining degrees of freedom.
\end{enumerate}
\end{frame}
\note{All scores are hypothetical teaching data. Solution to Problem 1. The sample mean is 70; the deviations are -10,-10,0,10,10 and sum to zero. The sum of squared deviations is 400. The sample variance is 400/(5-1)=100 points squared and the sample standard deviation is sqrt(100)=10 points. Dividing by n instead gives 80 points squared. Estimating the mean imposes one linear constraint on the deviations, leaving n-1=4 degrees of freedom. Under the stated i.i.d. finite-variance assumptions, the expected sum of squared deviations is (n-1) times the population variance. Bessel's correction makes the sample variance unbiased; it does not generally make the sample standard deviation unbiased.}

\ifshowsolutions
% SOLUTION 1 START
\begin{frame}{Solution 1: Mean, deviations, and sample dispersion}
\small
\answerpart{1}The sample mean is
\[
\overline y=\frac1n\sum_{i=1}^n y_i
=\frac{60+60+70+80+80}{5}=\boxed{70\text{ points}}.
\]
\begin{center}
\begin{tabular}{c|rrrrr|r}
Student $i$ & 1 & 2 & 3 & 4 & 5 & Sum \\ \hline
$y_i-\overline y$ & $-10$ & $-10$ & 0 & 10 & 10 & 0 \\
$(y_i-\overline y)^2$ & 100 & 100 & 0 & 100 & 100 & 400
\end{tabular}
\end{center}
\answerpart{2}Using the sum of squared deviations,
\[
s_Y^2=\frac{\sum_{i=1}^n(y_i-\overline y)^2}{n-1}
=\frac{400}{5-1}=\boxed{100\text{ points}^2}.
\]
The standard deviation is the square root of the variance:
\[
s_Y=\sqrt{s_Y^2}=\sqrt{100}=\boxed{10\text{ points}}.
\]
\end{frame}

\begin{frame}{Solution 1: Why the divisor is $n-1$ (continued)}
\small
\answerpart{3}Dividing by $n$ would give
\[
\widetilde s_Y^2=\frac1n\sum_{i=1}^n(y_i-\overline y)^2
=\frac{400}{5}=\boxed{80\text{ points}^2},
\]
compared with $s_Y^2=100$ points squared using $n-1$.
\begin{itemize}
\setlength{\itemsep}{7pt}
\item Once the mean is fitted, $\sum_i(y_i-\overline y)=0$. Knowing four of the five deviations determines the fifth: there are $n-1=4$ degrees of freedom.
\item For i.i.d. observations with finite population variance $\sigma_Y^2$,
\[
\E\!\left[\sum_{i=1}^n(Y_i-\overline Y)^2\right]=(n-1)\sigma_Y^2.
\]
Consequently, dividing by $n$ has expectation $\frac{n-1}{n}\sigma_Y^2$, whereas dividing by $n-1$ has expectation $\sigma_Y^2$.
\end{itemize}
\begin{block}{Bessel's correction}
The divisor $n-1$ makes the \emph{variance estimator} unbiased. Its square root $s_Y$ is not generally an unbiased estimator of $\sigma_Y$.
\end{block}
\end{frame}
% SOLUTION 1 END
\fi

\begin{frame}{Problem 2: How precise is the average test score?}
Test scores $Y_1,\ldots,Y_n$ are i.i.d. with population standard deviation $\sigma_Y=10$ points. Initially, $n=100$.
\begin{enumerate}
\setlength{\itemsep}{10pt}
\item Find $\Var(\overline Y)$ and the theoretical standard deviation of $\overline Y$.
\item Now suppose $\sigma_Y$ is unknown and the sample standard deviation is $s_Y=10$ points. Estimate the standard error of $\overline Y$. How does it differ from the standard deviation of individual scores?
\item Repeat both calculations for $n=400$. For this comparison, keep the population standard deviation and the supplied sample standard deviation at 10 points. What happens to precision when the sample size quadruples?
\end{enumerate}
\end{frame}
\note{All numerical quantities are hypothetical teaching data. Solution to Problem 2. For i.i.d. observations, Var(Ybar)=sigma squared/n. At n=100, the variance is 100/100=1 point squared and the theoretical SD is sqrt(1)=1 point. When sigma is unknown, SE(Ybar)=sqrt(s squared/n)=s/sqrt(n)=1 point. Individual scores have supplied SD 10 points, whereas the sampling distribution of the mean has SD 1 point. At n=400, the variance is .25 points squared, the theoretical SD is sqrt(.25)=.5 points, and the estimated SE is .5 points when s remains 10. Quadrupling n halves the SD and SE; actual sample standard deviations need not remain fixed across samples.}

\ifshowsolutions
% SOLUTION 2 START
\begin{frame}{Solution 2: Individual dispersion and uncertainty in an average}
\small
\answerpart{1}For i.i.d. scores, the variance of the mean and its square root are
\[
\Var(\overline Y)=\frac{\sigma_Y^2}{n}=\frac{10^2}{100}=1\text{ point}^2,
\qquad
\operatorname{SD}(\overline Y)=\sqrt{\frac{\sigma_Y^2}{n}}=\sqrt1=1\text{ point}.
\]
\answerpart{2}When $\sigma_Y$ is unknown, replace it with $s_Y$:
\[
SE(\overline Y)=\sqrt{\frac{s_Y^2}{n}}=\frac{s_Y}{\sqrt n}
=\frac{10}{\sqrt{100}}=\boxed{1\text{ point}}.
\]
The supplied individual-score SD is 10 points; the estimated uncertainty in the sample mean is 1 point.
\par\vspace{5pt}
\answerpart{3}With $n=400$, holding both supplied SDs at 10 points,
\begin{align*}
\Var(\overline Y)&=\frac{\sigma_Y^2}{n}=\frac{100}{400}=.25\text{ points}^2,\\
\operatorname{SD}(\overline Y)&=\sqrt{\frac{\sigma_Y^2}{n}}=\sqrt{.25}=.50\text{ points},\\
SE(\overline Y)&=\sqrt{\frac{s_Y^2}{n}}=\sqrt{\frac{100}{400}}=.50\text{ points}.
\end{align*}
Quadrupling $n$ quarters the variance and halves the SD and estimated SE.
\end{frame}
% SOLUTION 2 END
\fi

\section{Testing}

\begin{frame}{Problem 3: Is the average score different from 70?}
An i.i.d. sample of 100 students has mean score $\overline y=72.4$ and sample standard deviation $s_Y=10$ points. Use the large-sample normal approximation.
\begin{enumerate}
\setlength{\itemsep}{7pt}
\item State the null and alternative hypotheses for a two-sided test of whether the population mean is 70 points.
\item Calculate the estimated standard error, the observed $t$-statistic, and the two-sided p-value. You may use $\Phi(2.4)=.9918$, where $\Phi$ is the standard normal c.d.f.
\item Decide whether to reject at the 10\%, 5\%, and 1\% levels. Check your decisions using the corresponding two-sided cutoffs $1.645$, $1.96$, and $2.576$.
\item Interpret the p-value conditional on the null hypothesis. Does it give the probability that the null hypothesis is true?
\end{enumerate}
\end{frame}
\note{All sample statistics are hypothetical teaching data. Solution to Problem 3. The null is mu=70 and the alternative is mu not equal to 70. SE=s/sqrt(n)=10/10=1 point, and the observed t-statistic is (72.4-70)/1=2.4. Using the large-sample normal approximation, p=2[1-Phi(2.4)]=2(.0082)=.0164. Reject at 10 percent and 5 percent; fail to reject at 1 percent. The respective critical values are 1.645, 1.96, and 2.576. Under the null and the approximation, the probability of obtaining a test statistic at least as extreme in either direction is about 1.64 percent. This is not a posterior probability that the null is true. Do not use an exact Student t distribution for this exercise.}

\ifshowsolutions
% SOLUTION 3 START
\begin{frame}{Solution 3: Hypotheses, standard error, and the p-value}
\small
\answerpart{1}We test the population mean $\mu_Y$:
\[
H_0:\mu_Y=70\qquad\text{against}\qquad H_1:\mu_Y\ne70.
\]
\answerpart{2}The estimated standard error is the square root of the estimated sampling variance:
\[
SE(\overline Y)=\sqrt{\frac{s_Y^2}{n}}=\frac{s_Y}{\sqrt n}
=\frac{10}{\sqrt{100}}=1\text{ point}.
\]
\[
t^{\mathrm{obs}}=\frac{\overline y-\mu_Y^0}{SE(\overline Y)}
=\frac{72.4-70}{1}=\boxed{2.4}.
\]
Under $H_0$, the large-sample approximation gives $T\approx\mathcal N(0,1)$. Thus
\begin{align*}
\text{p-value}&=\Pr_{H_0}(|T|\geq|t^{\mathrm{obs}}|)\\
&\approx2\bigl[1-\Phi(|t^{\mathrm{obs}}|)\bigr]
=2[1-.9918]=\boxed{.0164}.
\end{align*}
Both tails count because $H_1$ allows means above or below 70.
\end{frame}

\begin{frame}{Solution 3: Decisions and the meaning of .0164 (continued)}
\small
\answerpart{3}Reject $H_0$ when the p-value is at most $\alpha$, or equivalently when $|t^{\mathrm{obs}}|$ reaches the two-sided critical value.
\begin{center}
\renewcommand{\arraystretch}{1.3}
\begin{tabular}{c|c|c|l}
Level $\alpha$ & Critical value & Compare $|t^{\mathrm{obs}}|=2.4$ & Decision \\ \hline
10\% & 1.645 & $2.4>1.645$ & Reject: $.0164<.10$ \\
5\% & 1.960 & $2.4>1.960$ & Reject: $.0164<.05$ \\
1\% & 2.576 & $2.4<2.576$ & Fail to reject: $.0164>.01$
\end{tabular}
\end{center}
\vspace{.15cm}
\answerpart{4}If the population mean were 70, the probability of a test statistic at least as far from zero as $2.4$, in either direction, would be approximately \textbf{1.64\%} under the large-sample approximation.
\begin{block}{What is being conditioned on?}
The calculation assumes $H_0$ to obtain a probability of data as extreme as the observed result. It does not calculate the probability that $H_0$ is true after observing the data.
\end{block}
\end{frame}
% SOLUTION 3 END
\fi

\begin{frame}{Problem 4: Does a larger sample change the evidence?}
Two independent surveys draw i.i.d. test scores from the same student population:
\begin{center}
\begin{tabular}{c|rrr}
Survey & Sample size $n$ & Mean $\overline y$ & Sample SD $s_Y$ \\ \hline
A & 100 & 71.5 & 10 \\
B & 400 & 71.5 & 10
\end{tabular}
\end{center}
In each survey, test $H_0:\mu_Y=70$ against $H_1:\mu_Y\ne70$ using the large-sample normal approximation.
\begin{enumerate}
\setlength{\itemsep}{8pt}
\item Calculate each survey's estimated standard error and observed $t$-statistic.
\item Find each two-sided p-value. Use $\Phi(1.5)=.9332$ and $\Phi(3)=.99865$.
\item Give the 5\% decision for each survey. Explain why the evidence differs even though both observed mean scores are 71.5.
\end{enumerate}
\end{frame}
\note{All survey summaries are hypothetical teaching data. Solution to Problem 4. The estimated SE is sqrt(s squared/n)=s/sqrt(n). Survey A has SE 1 and t=(71.5-70)/1=1.5; survey B has SE .5 and t=(71.5-70)/.5=3. The approximate two-sided p-values are 2(1-.9332)=.1336 and 2(1-.99865)=.0027. At 5 percent, fail to reject for A and reject for B. Both estimate a difference of 1.5 points from the null, but B has half the SE, so its standardized difference doubles. Sample sizes alone do not determine p-values when the observed means or SDs differ; this example holds those summaries fixed.}

\ifshowsolutions
% SOLUTION 4 START
\begin{frame}{Solution 4: The same difference measured with greater precision}
\small
\answerpart{1}For each survey, first take the square root of the estimated variance:
\[
SE(\overline Y)=\sqrt{\frac{s_Y^2}{n}}=\frac{s_Y}{\sqrt n},
\qquad t^{\mathrm{obs}}=\frac{\overline y-\mu_Y^0}{SE(\overline Y)}.
\]
\begin{center}
\renewcommand{\arraystretch}{1.25}
\begin{tabular}{c|c|c}
Survey & Estimated SE (points) & Observed $t$-statistic \\ \hline
A & $10/\sqrt{100}=1$ & $(71.5-70)/1=1.5$ \\
B & $10/\sqrt{400}=.5$ & $(71.5-70)/.5=3$
\end{tabular}
\end{center}
\answerpart{2}Using $\text{p-value}\approx2[1-\Phi(|t^{\mathrm{obs}}|)]$,
\[
p_A\approx2(1-.9332)=\boxed{.1336},\qquad
p_B\approx2(1-.99865)=\boxed{.0027}.
\]
\answerpart{3}At 5\%, A \textbf{fails to reject} ($.1336>.05$); B \textbf{rejects} ($.0027<.05$).
\par\vspace{5pt}
Both estimate a difference of 1.5 points. With the same supplied sample SD, quadrupling the sample size halves the SE, so the observed $t$-statistic doubles.
\end{frame}
% SOLUTION 4 END
\fi

\begin{frame}{Problem 5: Which testing mistakes can we make?}
Scores are i.i.d. normal with mean $\mu$ and known standard deviation $\sigma=10$. With $n=100$, test
\[
H_0:\mu=70 \qquad\text{against}\qquad H_1:\mu\neq70.
\]
Reject if $|Z|\geq1.96$, where $Z=(\overline Y-70)/1$.
\begin{enumerate}
\setlength{\itemsep}{6pt}
\item Find the sample means for which we do not reject. Compute the Type I error probability when $\mu=70$.
\item If the true mean is $\mu=73$, compute the Type II error probability $\beta$ and the power of this test.
\item In 1,000 repetitions, how many Type I errors do we expect when $\mu=70$? How many detections do we expect when $\mu=73$?
\end{enumerate}
\vspace{3pt}
Use $\Phi(1.96)=.9750$, $\Phi(-1.04)=.1492$, and $\Phi(-4.96)\approx0$.
\end{frame}
\note{Solution to Problem 5. All data are invented for teaching. Exact normality with known sigma makes this calculation exact apart from the rounded cutoff and normal probabilities. SE=10/sqrt(100)=sqrt(1)=1 point. Do not reject for 68.04<ybar<71.96. Alpha=2[1-Phi(1.96)]=.05. At true mu=73, beta=Phi((71.96-73)/1)-Phi((68.04-73)/1)=Phi(-1.04)-Phi(-4.96), approximately .1492. Power approximately .8508. Expected false rejections under mu=70 are 1000(.05)=50; expected detections under mu=73 are 1000(.8508)=850.8, approximately 851. These are two separate hypothetical repetition settings. The error probabilities condition on the true mean; they are not posterior probabilities that a hypothesis is true. Counts fluctuate across repetitions.}

\ifshowsolutions
% SOLUTION 5 START
\begin{frame}{Solution 5: Rejection region and Type I error}
\small
\answerpart{1}The sampling standard deviation of the mean is
\[
SE(\overline Y)=\frac{\sigma}{\sqrt n}
=\sqrt{\frac{\sigma^2}{n}}
=\sqrt{\frac{100}{100}}=\sqrt1=1\text{ point}.
\]
We do not reject when
\[
\left|\frac{\overline Y-70}{1}\right|<1.96
\quad\Longleftrightarrow\quad
\boxed{68.04<\overline Y<71.96}.
\]
When $\mu=70$, $Z\sim\mathcal N(0,1)$, so
\begin{align*}
\alpha&=\Pr(\text{reject }H_0\mid\mu=70)\\
&=\Pr(|Z|\geq1.96)=2[1-\Phi(1.96)]\\
&=2(1-.9750)=\boxed{.05}.
\end{align*}
A Type I error means rejecting the null when the null is true.
\end{frame}

\begin{frame}{Solution 5: Type II error and power (continued)}
\small
\answerpart{2}When the true mean is 73, $\overline Y\sim\mathcal N(73,1)$, with variance 1. Hence
\begin{align*}
\beta(73)&=\Pr(\text{do not reject }H_0\mid\mu=73)\\
&=\Pr(68.04<\overline Y<71.96\mid\mu=73)\\
&=\Phi\!\left(\frac{71.96-73}{1}\right)
-\Phi\!\left(\frac{68.04-73}{1}\right)\\
&=\Phi(-1.04)-\Phi(-4.96)\approx\boxed{.1492}.
\end{align*}
\[
\text{Power}(73)=1-\beta(73)\approx1-.1492=\boxed{.8508}.
\]
\answerpart{3}Expected counts equal repetitions times the relevant probability:
\[
\begin{array}{ll}
\mu=70: & 1{,}000\alpha=1{,}000(.05)=\boxed{50}\text{ Type I errors},\\[3pt]
\mu=73: & 1{,}000[1-\beta(73)]\approx850.8\approx\boxed{851}\text{ detections}.
\end{array}
\]
These are expectations in two different settings, not guaranteed counts.
\end{frame}
% SOLUTION 5 END
\fi

\section{Intervals}

\begin{frame}{Problem 6: Which mean scores remain plausible?}
A large i.i.d. sample of $n=100$ students gives $\overline y=72.4$ and sample standard deviation $s_Y=10$ points. Use the normal approximation.
\begin{enumerate}
\setlength{\itemsep}{9pt}
\item Calculate the 90\%, 95\%, and 99\% confidence intervals for $\mu$. Use critical values 1.645, 1.96, and 2.576, respectively.
\item Use the 95\% interval to test each null against a two-sided alternative at the 5\% level: $H_0:\mu=70$ and $H_0:\mu=73$. Explain the connection between the interval and the tests.
\item Interpret 95\% coverage. Across 200 independently repeated samples, how many 95\% intervals would you expect to contain the true mean?
\end{enumerate}
\end{frame}
\note{Solution to Problem 6. All data are invented for teaching. Estimated SE=s/sqrt(n)=10/sqrt(100)=sqrt(100/100)=1. The interval is ybar plus or minus z times estimated SE. 90 percent: [70.755,74.045]; 95 percent: [70.44,74.36]; 99 percent: [69.824,74.976]. At 5 percent reject mu=70 because it lies outside the 95 percent interval; do not reject mu=73 because it is inside. Equivalently t70=2.4 exceeds 1.96 and abs(t73)=.6 is below 1.96. The interval comprises null mean values not rejected by the corresponding two-sided test. The parameter is fixed; the interval changes across samples. In repeated sampling approximately 95 percent of intervals contain the true mean, so expected count in 200 repetitions is 190; no guarantee about a particular batch. Coverage is approximate because estimated SE and normal critical values are used.}

\ifshowsolutions
% SOLUTION 6 START
\begin{frame}{Solution 6: Confidence levels and interval width}
\small
\answerpart{1}First compute the estimated standard error:
\[
SE(\overline Y)=\frac{s_Y}{\sqrt n}
=\sqrt{\frac{s_Y^2}{n}}
=\sqrt{\frac{100}{100}}=\sqrt1=1\text{ point}.
\]
The normal approximate interval is
\[
CI_{1-\alpha}=\left[\overline y-z_{1-\alpha/2}SE(\overline Y),\;
\overline y+z_{1-\alpha/2}SE(\overline Y)\right].
\]
\begin{center}
\begin{tabular}{ccl}
Confidence & Calculation & Interval (points)\\ \hline
90\% & $72.4\pm1.645(1)$ & $\boxed{[70.755,\;74.045]}$\\[7pt]
95\% & $72.4\pm1.96(1)$ & $\boxed{[70.44,\;74.36]}$\\[7pt]
99\% & $72.4\pm2.576(1)$ & $\boxed{[69.824,\;74.976]}$
\end{tabular}
\end{center}
Greater coverage requires a wider interval for the same data.
\end{frame}

\begin{frame}{Solution 6: Test inversion and repeated coverage (continued)}
\small
\answerpart{2}A two-sided 5\% test rejects $\mu_0$ when
\[
\left|\frac{\overline y-\mu_0}{SE(\overline Y)}\right|>1.96
\quad\Longleftrightarrow\quad
\mu_0\notin[\overline y-1.96SE,\;\overline y+1.96SE].
\]
Thus the 95\% interval contains the null values that are not rejected:
\[
\begin{array}{ccl}
\mu_0=70:& t=(72.4-70)/1=2.4, & \boxed{\text{reject }H_0},\\[4pt]
\mu_0=73:& t=(72.4-73)/1=-.6, & \boxed{\text{do not reject }H_0}.
\end{array}
\]
\answerpart{3}For intervals constructed by this method,
\[
\Pr\!\left(\mu\in[\overline Y-1.96SE(\overline Y),\;
\overline Y+1.96SE(\overline Y)]\right)\approx.95.
\]
The mean $\mu$ is fixed; the endpoints vary across samples. Approximately 95\% of intervals cover $\mu$ in repeated sampling. Among 200 repetitions, the expected number covering $\mu$ is
\[
200(.95)=\boxed{190}.
\]
\end{frame}
% SOLUTION 6 END
\fi

\begin{frame}{Problem 7: How many students give the desired precision?}
Plan an i.i.d. survey of scores using a preliminary standard deviation of $\sigma=10$ points. Use a 95\% normal approximate confidence interval.
\begin{enumerate}
\setlength{\itemsep}{10pt}
\item Derive the minimum sample size needed for an interval half-width of at most 1 point. Round up to a whole student.
\item How large must the sample be for a half-width of at most .5 points? Explain the relationship between precision and sample size.
\item If the planning standard deviation is 20 points instead, how large a sample is needed for a half-width of at most 1 point? Explain the role of the planning assumption.
\end{enumerate}
\end{frame}
\note{Solution to Problem 7. All data are invented for teaching. Planning uses m=1.96 sigma/sqrt(n), giving n at least (1.96 sigma/m) squared. With sigma10,m1:384.16 so385. With m.5:1536.64 so1537. Halving half-width requires four times sample size before rounding. With sigma20,m1 same1536.64 and1537; doubling SD quadruples required n. Sigma is a preliminary planning value. If realized s differs from planned sigma, realized CI width changes, so n385 does not promise a half-width of exactly1.}

\ifshowsolutions
% SOLUTION 7 START
\begin{frame}{Solution 7: Sample size for a target half-width}
\small
\answerpart{1}The planning standard error and interval half-width are
\[
SE(\overline Y)=\sqrt{\frac{\sigma^2}{n}}=\frac{\sigma}{\sqrt n},
\qquad m_n=1.96\frac{\sigma}{\sqrt n}.
\]
For $m_n\leq m$, solve and round up:
\[
n\geq\left(\frac{1.96\sigma}{m}\right)^2
\quad\Longrightarrow\quad
n_{\min}=\left\lceil\left(\frac{1.96\sigma}{m}\right)^2\right\rceil.
\]
\[
\sigma=10,\;m=1:
\quad n_{\min}=\lceil(19.6)^2\rceil=\lceil384.16\rceil=\boxed{385}.
\]
\answerpart{2}For $\sigma=10$ and $m=.5$,
\[
n_{\min}=\lceil(1.96\cdot10/.5)^2\rceil
=\lceil1536.64\rceil=\boxed{1537}.
\]
Halving the target half-width multiplies required $n$ by four before rounding.
\vspace{6pt}

\answerpart{3}For $\sigma=20$ and $m=1$, $n_{\min}=\lceil(39.2)^2\rceil=\boxed{1537}$.
Doubling $\sigma$ also multiplies required $n$ by four. Realized precision depends on how well the planning value describes the population.
\end{frame}
% SOLUTION 7 END
\fi

\section{Groups}

\begin{frame}{Problem 8: Do two groups have the same average score?}
Two independent i.i.d. samples give the following summaries. Use the large-sample normal approximation and let $\delta=\mu_A-\mu_B$.
\begin{center}
\begin{tabular}{c|rrr}
Group & $n$ & $\overline y$ & $s$\\ \hline
$A$ & 100 & 75 & 12\\
$B$ & 144 & 71 & 12
\end{tabular}
\end{center}
\begin{enumerate}
\setlength{\itemsep}{5pt}
\item Estimate $\delta$ and calculate its standard error.
\item Test $H_0:\delta=0$ against $H_1:\delta\neq0$. Find the test statistic and p-value. Decide at the 5\% and 1\% levels.
\item Construct a 95\% confidence interval for $\delta$.
\item Test $H_0:\delta=3$ against $H_1:\delta\neq3$ at 5\%. Check your decision using the confidence interval.
\end{enumerate}
Use $\Phi(2.561)\approx.99478$ and $\Phi(.640)\approx.7389$.
\end{frame}
\note{Solution to Problem 8. All data are invented for teaching. Delta estimate75-71=4 points. Independence means covariance of sample means is zero. Var(ybarA-ybarB)=sigmaA squared/nA+sigmaB squared/nB, estimated by144/100+144/144=2.44. SE=sqrt(2.44)=1.562049935. Do not divide these terms by sample sizes again. Constant d0 does not change variance. For d0=0,t=4/SE=2.560737599, p approximately.01044; reject5 percent but do not reject1 percent. CI=4 plus or minus1.96SE=[.938382127,7.061617873], rounded[.938,7.062]. For d0=3,t=1/SE=.6401844,p approximately.5222 using supplied normal probability. Do not reject5 percent;3 lies in the95 percentCI. No assignment mechanism is given, so the observed difference is not automatically a causal effect.}

\ifshowsolutions
% SOLUTION 8 START
\begin{frame}{Solution 8: Difference, standard error, and a zero-null test}
\small
\answerpart{1}The estimated difference is $\widehat\delta=\overline y_A-\overline y_B=75-71=\boxed4$ points.
Independence makes the covariance of the sample means zero:
\[
\Var(\overline Y_A-\overline Y_B)
=\frac{\sigma_A^2}{n_A}+\frac{\sigma_B^2}{n_B}.
\]
Estimate the variance, then take its square root:
\[
SE(\widehat\delta)
=\sqrt{\frac{s_A^2}{n_A}+\frac{s_B^2}{n_B}}
=\sqrt{\frac{144}{100}+\frac{144}{144}}
=\sqrt{2.44}\approx\boxed{1.562}\text{ points}.
\]
\answerpart{2}For a hypothesized difference $d_0$,
\[
t=\frac{\widehat\delta-d_0}{SE(\widehat\delta)},
\qquad p=2[1-\Phi(|t|)].
\]
With $d_0=0$, $t=4/\sqrt{2.44}\approx2.561$, and
\[
p\approx2(1-.99478)=\boxed{.01044}.
\]
Reject at 5\% because $.01044<.05$; do not reject at 1\% because $.01044>.01$.
\end{frame}

\begin{frame}{Solution 8: An interval and a nonzero null (continued)}
\small
\answerpart{3}The 95\% normal approximate confidence interval is
\begin{align*}
CI_{.95}(\delta)
&=[\widehat\delta-1.96SE(\widehat\delta),\;
\widehat\delta+1.96SE(\widehat\delta)]\\
&=[4-1.96\sqrt{2.44},\;4+1.96\sqrt{2.44}]\\
&\approx\boxed{[.938,\;7.062]}\text{ points}.
\end{align*}
\answerpart{4}Changing the null value to $d_0=3$ changes the numerator:
\[
t=\frac{\widehat\delta-3}{SE(\widehat\delta)}
=\frac{4-3}{\sqrt{2.44}}\approx.640.
\]
\[
p=2[1-\Phi(|t|)]\approx2(1-.7389)=\boxed{.5222}.
\]
Do not reject $\delta=3$ at 5\%. Consistently, 3 lies inside the 95\% interval.
\vspace{7pt}

Subtracting a constant $d_0$ does not change the variance:
$\Var(\widehat\delta-d_0)=\Var(\widehat\delta)$. The interval estimates $\delta$ and remains centered at $\widehat\delta=4$.
\end{frame}
% SOLUTION 8 END
\fi

\begin{frame}{Problem 9: Does randomly assigned tutoring improve scores?}
Students are individually randomized into independent tutoring ($T$) and control ($C$) groups, with 100 students in each. Treat observations as i.i.d. within each group. There is full take-up, no spillover, and no differential attrition.
\begin{center}
\begin{tabular}{lcc}
\toprule
 & Tutoring & Control \\
\midrule
Posttest mean (SD) & $74\ (10)$ & $70\ (10)$ \\
Baseline mean (SD) & $69.5\ (10)$ & $70\ (10)$ \\
\bottomrule
\end{tabular}
\end{center}
Use large-sample normal inference: $z_{.975}=1.96$, $\Phi(2.828)\approx.99766$, and $\Phi(.3536)\approx.6382$.
\begin{enumerate}
\setlength{\itemsep}{4pt}
\item Estimate the posttest mean difference and its standard error.
\item Test no posttest effect at 5\%. Find the two-sided $p$-value and 95\% confidence interval.
\item Test equality of baseline means at 5\%. What does the result establish?
\item Under the stated design, how can we interpret the posttest difference?
\end{enumerate}
\end{frame}
\note{Solution to Problem 9. These data are hypothetical teaching data. The posttest difference is 74-70=4 points. For independent groups the estimated variance is sT squared/nT+sC squared/nC=100/100+100/100=2; SE=sqrt(2)=1.414 points. For the two-sided null difference zero, t=4/sqrt(2)=2.828 and the large-sample normal p-value is about .00468. Reject at 5 percent. The 95 percent confidence interval is 4 plus or minus 1.96 sqrt(2), approximately [1.228,6.772]. The baseline difference is -.5 and uses the same SE, giving t=-.3536 and p about .7236 using the supplied normal probability. Fail to reject baseline equality; this neither establishes exact equality nor proves randomization. Individual random assignment with full take-up, no spillovers, and no differential attrition permits interpreting the posttest difference as an estimate of the average causal effect for the study population. A baseline balance test cannot replace the design evidence.}

\ifshowsolutions
% SOLUTION 9 START
\begin{frame}{Solution 9: Estimating and testing the tutoring effect}
\small
\answerpart{1}For the observed posttest difference $\widehat\Delta=\overline y_T-\overline y_C$,
\[
\widehat\Delta=74-70=4\text{ points},
\qquad
\widehat{\Var}(\widehat\Delta)=\frac{s_T^2}{n_T}+\frac{s_C^2}{n_C}.
\]
The groups are independent, so
\[
SE(\widehat\Delta)=\sqrt{\frac{s_T^2}{n_T}+\frac{s_C^2}{n_C}}
=\sqrt{\frac{10^2}{100}+\frac{10^2}{100}}=\sqrt{2}\approx1.414\text{ points}.
\]
\answerpart{2}Test $H_0:\Delta=0$ against $H_1:\Delta\ne0$:
\[
t=\frac{\widehat\Delta-0}{SE(\widehat\Delta)}=\frac4{\sqrt2}\approx2.828,
\qquad p\approx2[1-\Phi(|t|)]\approx.0047.
\]
Since $|t|>1.96$ and $p<.05$, reject $H_0$. The 95\% confidence interval is
\[
\widehat\Delta\pm1.96\,SE(\widehat\Delta)
=4\pm1.96\sqrt2\approx\boxed{[1.228,\ 6.772]}\text{ points}.
\]
\end{frame}

\begin{frame}{Solution 9: Baseline balance and a causal interpretation}
\small
\answerpart{3}Let $\Delta_B$ denote the population baseline mean difference. For $H_0:\Delta_B=0$ against $H_1:\Delta_B\ne0$,
\[
\widehat\Delta_B=69.5-70=-.5,
\qquad
SE(\widehat\Delta_B)=\sqrt{\frac{10^2}{100}+\frac{10^2}{100}}=\sqrt2.
\]
\[
t_B=\frac{\widehat\Delta_B}{SE(\widehat\Delta_B)}=\frac{-.5}{\sqrt2}\approx-.3536,
\qquad
p_B\approx2[1-\Phi(.3536)]\approx2(1-.6382)=.7236.
\]
Fail to reject at 5\%. The sample gives little evidence of a baseline mean difference; this does not establish exact equality or prove that assignment was random.
\vspace{.3cm}

\answerpart{4}Random assignment, full take-up, no spillovers, and no differential attrition allow the posttest difference to estimate the average causal effect of tutoring for the study population: an estimated increase of 4 points.

\vspace{.2cm}
The design supports this interpretation. The baseline test is a diagnostic, not a proof of the design.
\end{frame}
% SOLUTION 9 END
\fi

\section{Data}

\begin{frame}{Problem 10: What do the school score summaries describe?}
Nine ordered school scores are
\[
60,\quad65,\quad70,\quad70,\quad75,\quad80,\quad80,\quad85,\quad90.
\]
\begin{enumerate}
\setlength{\itemsep}{8pt}
\item Find the sample mean, median, minimum, and maximum.
\item Define each quartile as the median of its half of the ordered data, excluding the overall median. Find $Q_1$, $Q_3$, and the interquartile range.
\item Compute the sample variance and standard deviation. Show the denominator used in the variance.
\item A \emph{separate, larger district sample} reports skewness $.37$ and ordinary kurtosis $3.15$ (not excess kurtosis). Interpret these summaries using the normal benchmark.
\end{enumerate}
\end{frame}
\note{Solution to Problem 10. These scores are hypothetical teaching data. The sum is 675, so the mean is 75. The fifth score, the median, is 75; the minimum is 60 and maximum 90. Excluding the overall median, the lower half is 60,65,70,70, so Q1=(65+70)/2=67.5. The upper half is 80,80,85,90, so Q3=(80+85)/2=82.5. IQR=15. The squared deviations from 75 sum to 750. The sample variance is 750/(9-1)=93.75 points squared and SD=sqrt(93.75)=9.682 points. The skewness and kurtosis are supplied for a different, larger sample and must not be presented as calculated from these nine scores. Positive skewness .37 indicates modest right asymmetry. Ordinary kurtosis 3.15 is slightly above the normal benchmark of 3, indicating a slightly larger standardized fourth moment. These summaries alone do not specify exact tail probabilities or establish a normal distribution.}

\ifshowsolutions
% SOLUTION 10 START
\begin{frame}{Solution 10: Location, quartiles, and spread}
\small
\answerpart{1}The sample mean is
\[
\overline y=\frac1n\sum_{i=1}^n y_i=\frac{675}{9}=75\text{ points}.
\]
The median is the fifth ordered score, 75. The minimum is 60 and the maximum is 90.
\vspace{.25cm}

\answerpart{2}Exclude the overall median before splitting the observations:
\[
\underbrace{60,\ 65,\ 70,\ 70}_{\text{lower half}},\qquad
75,\qquad
\underbrace{80,\ 80,\ 85,\ 90}_{\text{upper half}}.
\]
\[
Q_1=\frac{65+70}{2}=67.5,
\qquad Q_3=\frac{80+85}{2}=82.5.
\]
\[
\operatorname{IQR}=Q_3-Q_1=82.5-67.5=\boxed{15\text{ points}}.
\]
\end{frame}

\begin{frame}{Solution 10: Variance and distribution shape (continued)}
\small
\answerpart{3}The sample variance uses the divisor $n-1$:
\[
\begin{aligned}
\sum_{i=1}^9(y_i-\overline y)^2
&=225+100+25+25+0+25+25+100+225=750,\\
s_y^2&=\frac{\sum_{i=1}^n(y_i-\overline y)^2}{n-1}
=\frac{750}{8}=93.75\text{ points}^2,\\
s_y&=\sqrt{s_y^2}=\sqrt{93.75}\approx9.682\text{ points}.
\end{aligned}
\]
\answerpart{4}For the \emph{separate, larger sample}, skewness $.37>0$ indicates modest asymmetry toward the right.

\vspace{.2cm}
Ordinary kurtosis $3.15$ is slightly above the normal benchmark of 3: the standardized fourth moment is slightly larger.

\vspace{.2cm}
These summaries do not determine exact tail probabilities. They are not the skewness and kurtosis of the nine listed scores.
\end{frame}
% SOLUTION 10 END
\fi

\begin{frame}{Problem 11: Do study hours and test scores move together?}
Five students provide paired observations of weekly study hours $X$ and test scores $Y$:
\begin{center}
\begin{tabular}{c|rrrrr}
Student & 1 & 2 & 3 & 4 & 5 \\
\hline
Hours $x_i$ & 1 & 2 & 3 & 4 & 5 \\
Score $y_i$ & 60 & 70 & 65 & 80 & 75
\end{tabular}
\end{center}
\begin{enumerate}
\setlength{\itemsep}{7pt}
\item Find $\overline x$ and $\overline y$. Tabulate the centered observations and their cross-products.
\item Compute the sample covariance and state its units.
\item Compute both sample standard deviations and the sample correlation. What association does it describe?
\item Convert hours to minutes, $X^*=60X$. Find the covariance with $Y$ and the correlation. Explain the change.
\end{enumerate}
\end{frame}
\note{Solution to Problem 11. These pairs are hypothetical teaching data. Means are 3 hours and 70 points. Centered hours are -2,-1,0,1,2 and centered scores -10,0,-5,10,5, yielding cross-products 20,0,0,10,10 with sum 40. Sample covariance is 40/(5-1)=10 hour-points. The hours squared deviations sum to 10; scores squared deviations sum to 250. Sample SDs are sqrt(10/4)=sqrt(2.5) and sqrt(250/4)=sqrt(62.5); their product is 12.5. Correlation is 10/12.5=.8, a positive linear association, which does not by itself establish causality. Converting hours to minutes multiplies sample covariance by 60, giving 600 minute-points, and multiplies the X standard deviation by 60. Correlation remains .8 because the positive scale factor cancels.}

\ifshowsolutions
% SOLUTION 11 START
\begin{frame}{Solution 11: Centered observations and sample covariance}
\small
\answerpart{1}The sample means are
\[
\overline x=\frac15\sum_{i=1}^5x_i=\frac{15}{5}=3\text{ hours},
\qquad
\overline y=\frac15\sum_{i=1}^5y_i=\frac{350}{5}=70\text{ points}.
\]
\begin{center}
\begin{tabular}{c|rrrrr|r}
Student & 1 & 2 & 3 & 4 & 5 & Sum \\
\hline
$x_i-\overline x$ & $-2$ & $-1$ & $0$ & $1$ & $2$ & $0$ \\
$y_i-\overline y$ & $-10$ & $0$ & $-5$ & $10$ & $5$ & $0$ \\
$(x_i-\overline x)(y_i-\overline y)$ & $20$ & $0$ & $0$ & $10$ & $10$ & $40$
\end{tabular}
\end{center}
\answerpart{2}The sample covariance is
\[
s_{XY}=\frac1{n-1}\sum_{i=1}^n(x_i-\overline x)(y_i-\overline y)
=\frac{40}{5-1}=\boxed{10\text{ hour-points}}.
\]
Observations with higher study hours tend to have higher scores in this sample.
\end{frame}

\begin{frame}{Solution 11: Correlation and a change of units (continued)}
\small
\answerpart{3}For either variable $V$, $s_V^2=\sum_i(v_i-\overline v)^2/(n-1)$. Thus
\[
\begin{aligned}
s_X&=\sqrt{\frac{4+1+0+1+4}{4}}=\sqrt{2.5}\approx1.581\text{ hours},\\
s_Y&=\sqrt{\frac{100+0+25+100+25}{4}}=\sqrt{62.5}\approx7.906\text{ points}.
\end{aligned}
\]
\[
r_{XY}=\frac{s_{XY}}{s_Xs_Y}=\frac{10}{\sqrt{2.5}\sqrt{62.5}}
=\frac{10}{12.5}=\boxed{.8}.
\]
This is a positive linear association; correlation alone does not establish a causal effect.
\vspace{.2cm}

\answerpart{4}A positive change of units gives
\[
s_{X^*Y}=60s_{XY}=600\text{ minute-points},\qquad s_{X^*}=60s_X,
\]
\[
r_{X^*Y}=\frac{60s_{XY}}{(60s_X)s_Y}=r_{XY}=.8.
\]
Covariance depends on units; the common scale factor cancels in correlation.
\end{frame}
% SOLUTION 11 END
\fi

\begin{frame}{Problem 12: Are scores different across disadvantaged districts?}
Districts are classified by their share of students receiving subsidized meals. The samples are independent, with i.i.d. districts within each group:
\begin{center}
\begin{tabular}{lccc}
\toprule
Share receiving subsidized meals & $n$ & Mean $\overline y$ & SD $s$ \\
\midrule
Low ($L$) & 200 & 780 & 50 \\
High ($H$) & 200 & 720 & 50 \\
\bottomrule
\end{tabular}
\end{center}
Use large-sample normal inference, $z_{.975}=1.96$ and $\Phi(3.291)\approx.9995$.
\begin{enumerate}
\setlength{\itemsep}{6pt}
\item Estimate $\Delta=\mu_L-\mu_H$ and its standard error.
\item Test $H_0:\Delta=0$ against $H_1:\Delta\ne0$ at 5\%. Bound the two-sided $p$-value using the supplied normal probability.
\item Construct a 95\% confidence interval for the mean difference.
\item Does this estimate the causal effect of subsidized meals on test scores? Explain.
\end{enumerate}
\end{frame}
\note{Solution to Problem 12. These are hypothetical illustrative district samples, not the empirical California dataset from the lecture. Both supplied 50s are sample standard deviations. Estimate gap 780-720=60 points. Independent groups give SE=sqrt(sL squared/nL+sH squared/nH)=sqrt(2500/200+2500/200)=sqrt(25)=5 points. Under H0 difference zero the test statistic is 60/5=12. For a two-sided large-sample normal test, p=2[1-Phi(12)]=2Phi(-12), which is positive and below .001 because 12 exceeds 3.291. It is approximately 3.55 times 10 to the power -33, but students need not calculate this. Reject H0 at 5 percent. The confidence interval is 60 plus or minus 1.96(5), [50.2,69.8] points. The samples compare observational groups. Income, family background, school resources and other confounders may differ; this is an association, not an identified causal effect of meals.}

\ifshowsolutions
% SOLUTION 12 START
\begin{frame}{Solution 12: The gap and its statistical significance}
\small
\answerpart{1}The estimated mean difference is
\[
\widehat\Delta=\overline y_L-\overline y_H=780-720=60\text{ points}.
\]
For independent groups,
\[
\begin{aligned}
SE(\widehat\Delta)&=\sqrt{\frac{s_L^2}{n_L}+\frac{s_H^2}{n_H}}\\
&=\sqrt{\frac{50^2}{200}+\frac{50^2}{200}}
=\sqrt{12.5+12.5}=\boxed{5\text{ points}}.
\end{aligned}
\]
\answerpart{2}Under $H_0:\Delta=0$,
\[
t=\frac{\widehat\Delta-0}{SE(\widehat\Delta)}=\frac{60}{5}=12.
\]
Since $|t|>3.291$, the two-sided normal $p$-value satisfies
\[
p\approx2[1-\Phi(12)]<2[1-\Phi(3.291)]\approx.001.
\]
Reject $H_0$ at 5\%. The $p$-value is very small; it is not exactly zero.
\end{frame}

\begin{frame}{Solution 12: A confidence interval and its interpretation}
\small
\answerpart{3}The large-sample 95\% confidence interval is
\[
\begin{aligned}
\widehat\Delta\pm1.96\,SE(\widehat\Delta)
&=60\pm1.96(5)\\
&=60\pm9.8=\boxed{[50.2,\ 69.8]}\text{ points}.
\end{aligned}
\]
The entire interval is positive, agreeing with rejection of the zero-difference null in the two-sided 5\% test.
\vspace{.35cm}

\answerpart{4}The low-share districts have higher mean scores in these observational data. The difference does not identify a causal effect of subsidized meals.

\vspace{.2cm}
Districts may differ in income, family background, school resources, and other determinants of scores. Statistical significance does not remove these confounding differences.
\end{frame}
% SOLUTION 12 END
\fi

\end{document}
